\par\noindent\begin{minipage}{\linewidth}
\small
\setlength{\tabcolsep}{4pt}
\renewcommand{\arraystretch}{1.13}
\captionof{table}{Corpus y muestras de análisis.}\label{tab:T01}
\begin{tabular}{@{}>{\raggedright\arraybackslash}p{2.9cm}>{\raggedright\arraybackslash}p{5.8cm}>{\raggedright\arraybackslash}p{6.45cm}@{}}
\toprule
\textbf{Muestra} & \textbf{Artículos y representación} & \textbf{Comparación y alcance} \\
\midrule
Corpus original & 500.000 identificadores: 400.000 base + 100.000 de refuerzo. Diez modelos; título y resumen. & 26 áreas, cinco períodos. Los resúmenes globales usan la base; el corpus completo no es proporcional a la población. \\
\addlinespace[7pt]
Tres entradas de texto & 52.000 identificadores; 400 por área y período, 2.000 por área. Título, resumen y ambos. & Mismos artículos y candidatos entre modelos. Incluye el piloto de 26.000 artículos. \\
\addlinespace[7pt]
Fragmento común & Otra muestra distinta de 52.000 identificadores; 400 por área y período. & Contenido habitual y común emparejado; tokenización propia de cada modelo. Distinta de la muestra de tres entradas. \\
\addlinespace[7pt]
Grupos de igual tamaño & 256 artículos por centro; 26 áreas, 217 especialidades o 26 centros de especialidades. & Dentro de grupos: 26 áreas y las mismas 183 especialidades con 128, 256 y 512 artículos. Objetos distintos de los centros. \\
\addlinespace[7pt]
Búsquedas emparejadas & 217 especialidades; 50 consultas fijas por cada una (10.850 en total). Diez selecciones de candidatos. & 256 candidatos en cada ámbito, incluidas las 50 consultas. El conjunto del área está condicionado por esas consultas. \\
\addlinespace[7pt]
Geometría / parejas de áreas & 2.000 por área; veinte mitades de 1.000 y cinco selecciones de 2.000. & 325 parejas y dos propiedades; 26 condiciones guardadas. Las medidas alternativas tienen distinta cobertura. \\
\bottomrule
\end{tabular}
\end{minipage}\par
\par\smallskip{\small Las filas se solapan y no deben sumarse. Períodos: 2000--04, 2005--09, 2010--14, 2015--19 y 2020--24. Cada comparación usa los mismos identificadores entre modelos. En la búsqueda, cada artículo de consulta se excluye a sí mismo. Los 1.300 artículos de prueba técnica no forman otra muestra científica. Filtros y cuotas completos: Tabla S1.\par}
